We have presented a new multi-modal generative model, Multimodal Latent Diffusion (\gls{MLD}), to address the well-known coherence--quality tradeoff that is inherent in existing multi-modal \gls{VAE}-based models. \gls{MLD} uses a set of independently trained, uni-modal, deterministic autoencoders. Generative properties of our model stem from a masked diffusion process that operates on latent variables. We also developed a new multi-time training method to learn the conditional score network for multi-modal diffusion. An extensive experimental campaign on various real-life data-sets, provided compelling evidence on the effectiveness of \gls{MLD} for multi-modal generative modeling. In all scenarios, including cases with loosely correlated modalities and high-resolution datasets, \gls{MLD} consistently outperformed the alternatives from the state-of-the-art. 
% A limitation of our approach stems from the simple nature of encoder/decoder architectures. Concurrent work such as \citep{bao2023transformer} is more specialized, by focusing on complex, tailor-made encoder/decoder architectures. Such complexity might be necessary when moving to higher-resolution data. As for all generative models, also ours could be misused to produce misinformation. We believe however that the benefits of multi-modal generative models out-weight their potential mis-use.

% We have addressed the challenge of multi-modal generative modeling by proposing a novel method, Multimodal Latent Diffusion (\gls{MLD}). Our approach overcomes the coherence--quality tradeoff that is inherent in existing multi-modal \gls{VAE}-based models, whose limitations originate --- as we thoroughly discussed --- from problems related to latent variable collapse and information loss. To overcome these limitations, \gls{MLD} uses a set of independently trained, uni-modal, deterministic autoencoders. Generative properties of our model stem from a masked diffusion process that operates on latent variables. We also  developed a new multi-time training method to learn the conditional score network for multi-modal diffusion. An extensive experimental campaign on various real-life data-sets, provided compelling evidence on the effectiveness of \gls{MLD} for multi-modal generative modeling. In all scenarios, including cases with loosely correlated modalities, \gls{MLD} consistently outperformed the alternatives. A limitation of our approach stems from the simple nature of encoder/decoder architectures. Concurrent work such as \citep{bao2023transformer} is more specialized, by focusing on complex, tailor-made encoder/decoder architectures. Such complexity might be necessary when moving to higher-resolution data. As for all generative models, also ours could be misused to produce misinformation. We believe however that the benefits of multi-modal generative models out-weight their potential mis-use.
